% Use the existing owner-derived source mechanism; do not edit frozen originals.
\OLSADeclareDocumentTextCorrection{OLP-0009}{FA0009-ORDER-001}
{از اصل گسترش نتیجه می‌شود که اعضای یک مجموعه ترتیبی ندارند.}
{از اصل گسترش نتیجه می‌شود که خودِ مجموعه ترتیبی برای اعضایش تعیین نمی‌کند.}
\OLSADeclareDocumentTextCorrection{OLP-0009}{FA0009-NAME-002}
{وینر--کوراتفسکی}
{وینر--کوراتوفسکی}
\OLSADeclareDocumentTextCorrection{OLP-0009}{FA0009-COUNT-003}
{اگر $A$ دارای $n$ !!{element}s و $B$ دارای $m$ !!{element}s باشد،}
{اگر $A$ دارای $n$ !!{element} و $B$ دارای $m$ !!{element} باشد،}
\OLSADeclareDocumentTextCorrection{OLP-0009}{FA0009-COUNT-004}
{تعداد $m$ !!{element}s به صورت}
{تعداد $m$ !!{element} به صورت}
\OLSADeclareDocumentTextCorrection{OLP-0009}{FA0009-COUNT-005}
{$n\cdot m$ !!{element}s است.}
{$n\cdot m$ !!{element} است.}
\OLSADeclareDocumentTextCorrection{OLP-0009}{FA0009-COUNT-006}
{$n$ !!{element}s باشد، آنگاه $A^k$ دارای $n^k$ !!{element}s است.}
{$n$ !!{element} باشد، آنگاه $A^k$ دارای $n^k$ !!{element} است.}
\OLSADeclareDocumentTextCorrection{OLP-0009}{FA0009-FINITE-007}
{اگر $A$ مجموعه‌ای باشد، \emph{واژه‌ای} بر روی~$A$ هر دنباله‌ای از}
{اگر $A$ مجموعه‌ای باشد، \emph{واژه‌ای} بر روی~$A$ هر دنبالهٔ متناهی از}
\OLSADeclareDocumentTextCorrection{OLP-0009}{FA0009-WORD-008}
{پس مجموعهٔ
\emph{همهٔ} واژه‌ها بر روی~$A$ برابر است با}
{با این قراردادِ نمایشی، مجموعهٔ
\emph{همهٔ} واژه‌ها بر روی~$A$ را چنین می‌نویسیم\footnote{یادداشت ویراستاری: در این مثال، دنباله‌ها متناهی‌اند و طول، بخشی از مشخصات هر واژه است. فرمول زیر را باید اختصاری برای گردآوردن واژه‌های همهٔ طول‌ها دانست، نه تساویِ بی‌قیدوشرط میان کدهای مجموعه‌ایِ معرفی‌شده در بالا. برای الفبای دلخواه، برخوردِ کدها ممکن است رخ دهد: اگر مجموعهٔ تهی حرفی از الفبا باشد، کدِ واژهٔ تهی و کدِ واژهٔ تک‌حرفیِ آن یکسان است. همچنین اگر کدِ زوج مرتبِ دو حرف از الفبا نیز در الفبا باشد، کدِ واژهٔ تک‌حرفیِ آن زوج با کدِ واژهٔ دوحرفیِ متناظر یکسان است؛ آن دو حرف لازم نیست متمایز باشند. در تعریف دقیق، هر واژه را تابعی از یک قطعهٔ آغازینِ متناهی از اعداد طبیعی به الفبا می‌گیریم. دامنهٔ تابع طول را مشخص می‌کند؛ واژهٔ تهی تابعِ تهی است و با واژهٔ تک‌حرفی‌ای که حرف آن مجموعهٔ تهی است، تفاوت دارد. این همان قراردادِ طول‌دار برای رشته‌های متناهی است.}:}
\OLSADeclareSetup{OLP-0009}{\olsa_source_correction_install_document_text:}
